\documentclass[10pt,a4paper]{amsart}
\usepackage[margin=19mm]{geometry}
\usepackage{amsmath,amssymb,mathtools}
\usepackage[scr=rsfs,cal=euler]{mathalfa}
\usepackage{xfrac}
\usepackage[unicode,hidelinks,bookmarks=false]{hyperref}
\newcommand{\ie}{i.e.\ }
\newcommand{\cf}{cf.\ }
\newcommand{\resp}{respectively\ }
\newcommand{\Subsection}{\subsection}
\providecommand{\nobreakdash}{\nobreak\hbox{-}}
\setlength{\emergencystretch}{2em}
\multlinegap0pt
\usepackage{amssymb}
\usepackage{mathtools}
\newtheorem{lemma}{Lemma}
\DeclareMathOperator{\Pic}{Pic}
\DeclareMathOperator{\Sym}{Sym}
\DeclareMathOperator{\bC}{\mathbb{C}}
\newcommand{\hilb}[2][n]{#2^{[#1]}}
\newcommand{\sym}[2][n]{#2^{(#1)}}
\title{On Fujita's conjecture for a general hyperk\"ahler manifold in the standard series of examples}
\author{Alessandro Pilastro}
\date{Source: arXiv:2602.12439v1 (CC BY 4.0)}
\begin{document}
\maketitle

\noindent\emph{Source and licence.} On Fujita's conjecture for a general hyperk\"ahler manifold in the standard series of examples. Version arXiv:2602.12439v1 (CC BY 4.0), distributed by arXiv under the Creative Commons Attribution 4.0 International licence, http://creativecommons.org/licenses/by/4.0/, as recorded in the arXiv licence metadata for this exact version and shown on the abstract page https://arxiv.org/abs/2602.12439v1. Full licence notice: \texttt{licenses/arxiv-2602.12439-CC-BY-4.0.txt}.\par\smallskip

\noindent\textit{Editorial scope.} Counted unit: the article's lemma lemma\_bpf\_of\_L\_n\_on\_Hilb\_n (source0.tex lines 128--130). The article's complete original proof of it is delivered with it (source0.tex lines 131--136, one \texttt{proof} environment of 6 lines). No mathematics is written, repaired or re-derived: the statement and the proof are the article's own lines, in the article's own notation, and no retained line is substituted or re-flowed.  No further source-local context is needed: every cross-reference of every retained line resolves inside the two delivered spans.  Nothing was omitted from any retained span, so the package carries no omission record.  No retained line contains a citation, so the wrapper carries no bibliography and no bibliography entry.  No retained line contains a figure environment, an image-inclusion command or drawing code, so no image is delivered and the package declares no asset dependency.  Editorial formatting only, in addition: an AMS \texttt{amsart} preamble that restates the article's own declarations and macros used by the retained lines, namely the environments \texttt{lemma} on separate counters, together with the standard packages those lines need.  The retained environments print in this document as Lemma 1 in order of appearance, whereas the article prints the same items with the numbering of its own full document.  The complete native source of the article as distributed in the arXiv e-print for this exact version is bundled unchanged as \texttt{sources/194550/source0.tex}.
\begin{lemma}\label{lemma_bpf_of_L_n_on_Hilb_n}
Let $S$ be a smooth projective complex surface, and let $L\in\Pic(S)$ be base point free, then $L_n\in \Pic(\hilb{S})$ is also base point free.    
\end{lemma}

\begin{proof}
Let $\xi\in\hilb{S}$ and let $x_1,\dots,x_n$ be the points in the support of $\xi$ with multiplicities. We can find a section $s\in H^0(S,L)$ such that $s(x_i)\neq 0$ for all $i=1,\dots,n.$ This follows by induction on $n$. The base case is the definition of base point free line bundle. From the induction hypothesis, we have a section $\hat{s}$ such that $\hat{s}(x_i)\neq 0$ for $i=1,\dots,n-1$. If $\hat{s}(x_n)\neq 0$ then we are done. Suppose then that $\hat{s}(x_n)= 0$ and given $\tilde{s}$ a section non-vanishing at the point $x_n$, define $s=\alpha \hat{s}+\tilde{s}$. For a general $\alpha\in\bC$ the section $s$ satisfies the required property.

So, given $\xi\in\hilb{S}$ let $s$ be a section of $L$ non-vanishing at the points in the support of $\xi$. Consider now the section $\otimes_{i=1}^n\pi_i^*(s)\in H^0(S^n,L^{\boxtimes n})$ which is invariant under the action of $\Sym(n)$. It therefore descends to a section $\sym{s}\in H^0(\sym{S},\sym{L})$. Then $\sym{s}$ takes the value $\prod_{i=1}^ns(x_i)\neq0$ at the point $\sum_{i=1}^nx_i\in\sym{S}$. The pullback section $\mu^*(\sym{s})=\hilb{s}\in H^0(\hilb{S},L_n)$ does not vanish at the point $\xi\in\hilb{S}$ as required.

\end{proof}

\end{document}
